% Generated draft: requires mathematical and source review.
\documentclass[11pt]{article}
\usepackage[margin=25mm]{geometry}
\usepackage{fontspec}
\setmainfont{Latin Modern Roman}
\usepackage{amsmath,amssymb,mathtools,unicode-math}
\setmathfont{Latin Modern Math}
\usepackage{enumitem,needspace,xurl,hyperref}
\setlength{\parindent}{0pt}
\setlength{\parskip}{5pt}
\setlength{\emergencystretch}{3em}
\newcommand{\N}{\mathbb{N}}
\newcommand{\Z}{\mathbb{Z}}
\newcommand{\Q}{\mathbb{Q}}
\newcommand{\R}{\mathbb{R}}
\newcommand{\C}{\mathbb{C}}
\newcommand{\F}{\mathbb{F}}
\newcommand{\A}{\mathbb{A}}
\newcommand{\PP}{\mathbb{P}}
\newcommand{\E}{\mathbb{E}}
\newcommand{\eps}{\varepsilon}
\newcommand{\dd}{\,\mathrm{d}}
\begin{document}
\section*{UnsolvedMath ID: 1101204}
\textbf{Source problem code:} catalogue entry
\subsection{Definitions and mathematical statement}
For a connected orientable finite-type surface $S=S_{g,n}$ with $2g-2+n>0$, let $\mathcal T(S)$ be its Teichm\"uller space of marked complete finite-area hyperbolic structures, modulo the usual marking-preserving equivalence, and let $d_{\mathrm{WP}}$ be the path metric of the Weil--Petersson Riemannian metric, defined by the hyperbolic $L^2$ pairing on harmonic Beltrami differentials. Its complex dimension is $\xi(S)=3g-3+n$.

For $\xi(S)\geq1$, use the graph metric on the one-skeleton $\mathcal C(S)$ of the curve complex: vertices are isotopy classes of essential nonperipheral simple closed curves, with edges for disjoint representatives. For $S_{1,1}$ and $S_{0,4}$ use the usual modified adjacency of minimal intersection $1$ and $2$, respectively, giving the Farey graph. A geodesic metric space is Gromov hyperbolic if some $\delta\geq0$ makes all its geodesic triangles $\delta$-thin. Metric spaces $X,Y$ are quasi-isometric if there exist $f:X\to Y$, $A\geq1$, and $B\geq0$ such that
\[
A^{-1}d_X(x,x')-B\leq d_Y(f(x),f(x'))\leq A d_X(x,x')+B
\]
for all $x,x'\in X$, and every point of $Y$ is at distance at most $B$ from $f(X)$.

\textbf{Questions (Bowditch).} For which surfaces $S$ is $(\mathcal T(S),d_{\mathrm{WP}})$ Gromov hyperbolic? For which $S$ is it quasi-isometric to $\mathcal C(S)$? The original source does not specify the class of surfaces, so the conventional finite-type setting above is explicit rather than an assertion about all possible Teichm\"uller spaces.

The supplied summary reports negative answers to both questions when $\xi(S)\geq3$, and affirmative answers for $S_{1,1}$ and $S_{0,4}$ using the Farey graph. It attributes the negative hyperbolicity result to Brock--Farb and hyperbolicity of the curve complex to Masur--Minsky. Although called a clean dichotomy there, these cases omit $\xi(S)=2$ and do not settle that part from the supplied evidence. The summary identifies the pants graph as a quasi-isometric model for the Weil--Petersson metric, attributed to Brock.
\subsection{Short English statement}
Determine when Weil--Petersson Teichm\"uller space is Gromov hyperbolic and when its large-scale metric agrees with the curve graph. The supplied status summary distinguishes higher complexity from two Farey-graph cases but omits complexity two.
\subsection{Sources}
\begin{itemize}
\item M. Bestvina, \emph{Questions in Geometric Group Theory} (2004), Question 12.4, PDF page 20, attributed to Bowditch: \url{https://www.math.utah.edu/~bestvina/eprints/questions-updated.pdf}; manual transcription, accessed 2026-07-29.
\item Supplied research-summary references: Brock--Farb (2006), Weil--Petersson nonhyperbolicity for complex dimension at least three; Masur--Minsky (1999), hyperbolicity of the curve complex; Brock (2003), pants-graph model for the Weil--Petersson metric. No titles or theorem identifiers are supplied.
\item Local registry \texttt{data/unsolvedmath/problems.json}, permanent ID $1101204$, source code catalogue entry, \emph{Questions in Geometric Group Theory --- Q 12.4}. The record labels this solved and SOLVED-IN-LITERATURE; its source-list background retains current-status and rights-review qualifications.
\end{itemize}
\end{document}
